\begin{papercard}{gwet2008}
\cf{Stages} \stagetag{I}. A statistical paper that analyses the origin of the ``paradoxes of kappa'' (high agreement but low $\pi$/$\kappa$ when a category is rare or omnipresent), proposes the alternative \textbf{AC1} coefficient with a calibrated chance-agreement term, and gives variance estimators for the multi-rater $\pi$ and AC1 that do not assume rater independence, with a small Monte-Carlo study. Its core is mathematics; this card keeps the reasoning and the results.

\cf{Scope and notation} Inter-rater reliability for \emph{nominal} data. All coefficients share an overall agreement $p_a$ and differ in the chance term $p_{e}$: Cohen's $\kappa$ uses $\sum p_{Ak}p_{Bk}$, Scott's $\pi$ uses squared pooled proportions, the $G$-index uses $1/q$ (number of categories), and Fleiss' generalised $\pi$ extends $\pi$ to $r$ raters. Conger averages pairwise $\kappa$s \pc{gwet2008}{2}{All inter-rater reliability coefficients discussed in this paper have two components}.

\cf{The paradox} Table 3: 125 participants, 94.4\,\% agreement, yet $\pi=(0.9440-0.9456)/(1-0.9456)=-0.0288$, a negative value ``opposite to what our intuition would suggest'' (Cicchetti \& Feinstein's high agreement with low $\kappa$) \pc{gwet2008}{5}{illustrates one of the paradoxes noted by Cicchetti and Feinstein (1990) where high agreement is coupled with low}. A rater model with sensitivity, specificity and trait prevalence shows that overall agreement $P_a$ does \emph{not} depend on prevalence, whereas $\pi$ and $\kappa$ do: when a trait is very rare or very common, $\pi$ becomes negative or $\kappa$ falls to 0 regardless of rater quality, and both behave sensibly only near prevalence 0.5 \pc{gwet2008}{6}{whenever a trait is very rare or omnipresent, Scott}, \pc{gwet2008}{7}{A value of trait prevalence that is either close to 0 or close to 1 will considerably reduce the ability of k to reflect any extent of agreement between raters}. Diagnosis: ``the chance-agreement probability used in these statistics is ill-estimated when trait prevalence is in the neighbourhood of 0 or 1''.

\cf{AC1} Premises: chance agreement occurs only when at least one rater rates randomly, and only an unknown portion of ratings is random. Model: rater A rates randomly with propensity $u_A$ and otherwise correctly, so the ``true'' reliability is $\gamma=\dfrac{(\text{certain agreement})}{(\text{all but chance})}$. Chance agreement is approximated as $P(G\mid R)\cdot C$ with $P(G\mid R)=1/2$ and the normalised variance $C=4p_+(1-p_+)$, i.e.\ $p_e=2p_+(1-p_+)$ for two categories \pc{gwet2008}{9}{I suggest approximating P(R) with a normalized measure of randomness C defined by the ratio of the variance V(Xþ) of Xþ}. General form: $p_e=\frac{1}{q-1}\sum_k \hat p_k(1-\hat p_k)$. AC1 $=(p_a-p_e)/(1-p_e)$ is called ``first-order'' because it corrects for full agreement only (an AC2 is announced for weighted agreement). Chance correction is therefore \emph{proportional to the portion of ratings that could agree by chance} \pc{gwet2008}{10}{Unlike the k- and p-statistics, this agreement coefficient uses a chance-agreement probability that is calibrated}. At prevalence near 0.5 all indices behave alike.

\cf{Variance estimation} Linearisation and jackknife estimators for the generalised $\pi$, a pairwise-$\kappa$ generalisation and AC1, valid without independence between raters (existing variances assume independence). The AC1 variance is $\hat v(\hat\gamma_1)=\frac{1-f}{n}\frac{1}{n-1}\sum_i(\hat\gamma^*_{1|i}-\hat\gamma_1)^2$ with sampling fraction $f=n/N$ \pc{gwet2008}{14}{To derive its variance, I have used a linear approximation that includes all terms with a stochastic order of magnitude}. For Table 3: $\kappa$ and $\pi$ low and inconsistent with 94.4\,\% agreement, AC1 and $G$ consistent \pc{gwet2008}{16}{Because the percentage agreement pa equals 94.4\%, it appears that AC1 and G-index are more consistent with the observed extent of agreement}.

\cf{Monte-Carlo study} Two raters, $n=20,60,80,100$, prevalence $P_r=95\,\%$, propensity for random rating $u=5\,\%$ or $20\,\%$/$5\,\%$, 500 replicates; non-random ratings are correct and random ones correct with probability 1/2. Relative bias: AC1 $-0.8$ to $0.0\,\%$ ($-2.1$ to $-1.3\,\%$ with $u_A=20\,\%$), $G$-index $-3.6$ to $-12.1\,\%$, $\pi$ and $\kappa$ $-32.8$ to $-62.5\,\%$ \pc{gwet2008}{17}{It follows from Table 5 that the relative bias of the AC1 estimator}. AC1 had the smallest variance (0.07--0.33\,\%) vs.\ 2--15\,\% for $\pi$ and $\kappa$; the proposed variance estimators worked well, except that $\kappa$/$\pi$ estimators underestimate variance for small samples \pc{gwet2008}{19}{An examination of the Monte-Carlo expectation of the various variance estimators indicates that the proposed variance estimators for AC1 and G-index work very well}.

\cf{Conclusion and limits} $\kappa$ and $\pi$ treat all observed ratings as if they could agree by chance, which gives unpredictable results when the data suggest little chance agreement \pc{gwet2008}{19}{Kappa and pi evaluate the chance-agreement probability as if all observed ratings may yield an agreement by chance}. Limits: nominal two-category core model (extended to $q$ categories by formula), a simple random-rating model, a single simulated prevalence (95\,\%), and the ``true'' reliability defined within Gwet's own model, so the comparison is not model-free.

\cf{Use for this thesis} The standard citation for adding AC1 as a robustness statistic when privacy codes are rare (e.g.\ presence of a dark pattern in few apps) and $\kappa$/$\alpha$ collapse despite high percent agreement. The thesis should (1) report raw agreement, joint-presence counts, $\kappa$, $\alpha$ and AC1 together for skewed codes, (2) state \emph{in advance} which is primary, and (3) not choose the statistic after seeing the results. Krippendorff (2004, 2019) disputes AC1's interpretation (its zero point is at an arbitrary configuration and it yields higher numbers), so present AC1 as a sensitivity analysis rather than a replacement. AC1 is implemented in R (irrCAC), whose version must be recorded.
\end{papercard}
